\section{Experimental Methodology}
\label{sec:method}

\paragraph{Research questions.} RQ1: does SAT-SA emit the finding families declared
for controlled scenarios? RQ2: how does it compare with simple baselines, and what
does each worker contribute? RQ3: how does it behave on imperfect evidence? RQ4: how
sensitive is the peer rule to cohort size and dispersion? RQ5: what do durable
orchestration and review cost, and do interrupted runs recover? RQ6: does
verification detect alteration of the finalized record? EXT: does the pipeline run
on independent real data, and is its risk associated with an outcome available
there? Utility to supervisors was outside the evaluated scope because no
human-subject study was conducted. Table~\ref{tab:design} maps every question to
its experiment.

\begin{table}[t]
\caption{Research questions, experiments and independent units. ``Trials'' repeat a
measurement on the same unit and are not independent.}
\label{tab:design}
\centering
\scriptsize
\setlength{\tabcolsep}{3pt}
\begin{tabular}{@{}L{1.7cm}L{1.2cm}L{3.1cm}L{2.7cm}L{2.6cm}@{}}
\toprule
RQ & Exp. & Data (ground truth) & Baselines & Independent unit (n) \\
\midrule
RQ1 detection & C01 & catalog (construction labels) & -- & scenario (5) \\
RQ2 baselines & C01 & closure corpus (construction) & threshold, MAD, z, IQR, random & record (12) \\
RQ2 ranking & PR01 & generated populations (generator) & random, volume, fastest closure & population (20) \\
RQ2 ablation & A01 & catalog; populations & full worker set & scenario (5); population (5) \\
RQ4 robustness & R01b & perturbed catalog (declared contract) & -- & fixture (5); 245 conditions \\
RQ5 integrity & T01 & one workflow (declared outcome) & unsigned-field control & workflow (1); 14 mutations \\
RQ6 peer & P01 & engineered cohorts (design) & -- & designed cell (96) \\
RQ7 cost & O01 & catalog scenario & direct vs.\ graph vs.\ unreviewed & 1 machine; 30 paired trials \\
RQ7 recovery & O02 & catalog scenario (reference run) & uninterrupted run & 12 cells; 36 runs \\
EXT & X02b, X03 & UCI-498 (SLA flag) & random, volume, reassign., slowest & group of 1 organization (50) \\
\bottomrule
\end{tabular}
\end{table}


\paragraph{Data.} All controlled data are \emph{synthetic}. The \emph{controlled
catalog} contains 5 executable scenarios (of 10
catalog entries) authored with declared finding families and supervisory actions.
Its labels are \emph{construction labels}: they encode what a scenario was built to
contain, so agreement shows that the mechanism responds to its designed inputs, not
that it detects real misconduct. \emph{Generated populations} (PR01) are
20 independently seeded populations of 20 entities,
5 of them declared pathological by the generator.
\emph{Engineered peer cohorts} (P01) are 96 cells crossing cohort size,
dispersion and subject value. Controlled data were necessary because no public dataset
contains SOC workflow records with supervisory ground truth. The single \emph{external}
dataset is the public UCI incident-management event log (dataset 498, Creative Commons Attribution licence) of one
organization's IT service desk~\cite{amaral2018uci498,amaral2019completion}; it is IT
incident data, not SOC data. Two security datasets were assessed and rejected: GUIDE
has triage-graded SOC incidents but no acknowledgement, closure, investigation-step or
escalation records~\cite{freitas2024guide}, and CIC-IDS2017 contains network flows,
not workflow records~\cite{sharafaldin2018cicids}.

\paragraph{Baselines.} Baselines are deliberately simple because they provide
transparent reference points for assessing whether the additional analytical
machinery yields measurable benefit. For closure-time detection:
fixed threshold, MAD rule, z-score and interquartile-range (IQR) rules, and seeded
random. For prioritization: random order, critical-and-high alert volume, and a
closure-speed heuristic that ranks entities by \emph{fastest median closure time}.
Externally: random, incident volume, reassignment rate and slowest median
resolution; the last is close to the outcome's own construct and serves as an upper
reference, not a competitor.

\paragraph{Metrics and statistics.} Precision, recall and F1 against construction
labels; recall and normalized discounted cumulative gain
(NDCG)~\cite{jarvelin2002ndcg} within the top $k\in\{10,20,30\}\%$ of entities;
review volume needed to find all pathological entities; conformance with the declared
validation contract; wall-clock seconds and database statements; and Spearman $\rho$
against the share of incidents that missed their service-level agreement (SLA). We
report medians, 95th percentiles only when $n\ge 20$, and seeded percentile bootstrap
intervals~\cite{efron1993bootstrap} only when there are at least ten independent
units; comparisons are paired by seed or trial with counts of units where one arm was
higher, equal or lower. No $p$-values are computed and no significance is claimed; all
intervals are descriptive and uncorrected. Table~\ref{tab:design} distinguishes
repeated trials from independent units.

\paragraph{Evidence discipline.} Each experiment produces an experiment bundle
(manifest with code commit and a dirty-tree flag, configuration, raw results,
processed metrics) with SHA-256 hashes. The canonical bundles are frozen and
hash-verified; superseded bundles (for example X02, replaced by X02b after a
bookkeeping correction that left all analytical outputs identical) are kept but not
cited. Every number in this paper is rendered from the frozen bundles; no detector,
threshold or weight was changed on evaluation data.
